\begin{abstract}
High-resolution UAV imagery offers a scalable basis for nursery inventory, but it is poorly matched to the fixed visual context of compact vision--language models (VLMs). A single 5K frame may contain hundreds of small plants; whole-image downsampling removes instance detail, whereas exhaustive overlapping tiles repeatedly process the same regions and make inference cost grow with image size. Transfer across ornamental crops is also difficult because canopy geometry, spacing, object scale, and background composition change substantially. We introduce \emph{OrnaScout}, a budget-aware framework that allocates a fixed visual-computation budget to cross-crop nursery monitoring. Splitting precedes tiling to prevent source leakage. A supervised detector provides local evidence; the Budgeted Evidence Selector combines score-based exploitation with controlled exploration to retain the same number of tiles for every VLM. A schema-constrained decoder produces auditable records, and a validation-fitted gate routes unsupported records to human review. We study 418 relabeled Azalea and Magnolia images containing 29,842 plant instances and 14,630 source-safe tiles. Azalea has a median of 97 plants per image versus 57 for Magnolia, and its median normalized box area is 44.1\% smaller. In a source-level validation audit at an eight-tile budget, count-ranked selection retained 45.5\% of annotated evidence, compared with 35.2\% for random selection. Exploit--explore routing retained 43.9\%, reduced overlap redundancy by 22.7\% relative to combined-score top-$K$, and increased hard-negative retention from 0.7\% to 6.7\%. These results expose a trade-off among foreground recall, spatial diversity, and negative evidence, establishing an equal-budget protocol for cross-crop VLM and selective-risk evaluation.
\end{abstract}
